% Part: lambda-calculus
% Chapter: lambda-definability
% Section: fixpoints

\documentclass[../../../include/open-logic-section]{subfiles}

\begin{document}

\olfileid{lam}{ldf}{fp}

\olsection{Titik Tetep}

Upamakna kita arep netepake fungsi faktorial kanthi rekursi minangka
suku $\fn{Fac}$ kang nduweni sipat iki:
\[
\fn{Fac} \ident \lambd[n][\fn{IsZero}\, n\, \num 1 (\fn{Mult}\, n (\fn{Fac}(\fn{Pred} \, n)))]
\]
Tegese, faktorial saka $n$ yaiku $1$ yen $n = 0$, lan $n$ dikali
faktorial saka $n-1$ yen ora. Mesthi wae, kita ora bisa netepake suku
$\fn{Fac}$ kanthi cara iki amarga $\fn{Fac}$ dhewe katon ing sisih
tengen. Definisi rekursif kang ngrujuk marang awake dhewe kaya iki
dudu bagean saka kalkulus lambda. Netepake sawijining suku, kayata
\[
\fn{Mult} \ident \lambd[ab][a (\fn{Add}\, b) \num{0}]
\]
mung nggunakake suku kang wis ditetepake sadurunge ing sisih tengen,
kayata $\fn{Add}$. Kita bisa tansah ngilangi $\fn{Add}$ kanthi ngganteni
nganggo suku kang netepake. Kanthi mangkono, $\fn{Mult}$ dadi suku
lambda murni; yen $\fn{Add}$ dhewe ngemot suku kang wis ditetepake
(kaya $\fn{Add}'$), kita bisa nerusake proses iki nganti pungkasane
entuk suku lambda murni.

Nanging bab iki ora lumaku kanggo definisi rekursif kaya
$\fn{Fac}$ ing ndhuwur. Yen kedadeyan $\fn{Fac}$ ing sisih tengen
kita ganti nganggo definisi $\fn{Fac}$ dhewe, kita entuk:
\begin{align*}
  \fn{Fac} & \ident \lambd[n][\fn{IsZero}\, n\, \num{1}\,
    (\fn{Mult}\, n\,
      ((\lambd[n][\fn{IsZero}\, n\, \num{1}\,
        (\fn{Mult}\, n\, (\fn{Fac}(\fn{Pred}\, n)))])
       (\fn{Pred}\, n)))]
\end{align*}
lan $\fn{Fac}$ isih durung ilang saka sisih tengen.
Cetha yen proses iki dibaleni, definisine dadi saya dawa
lan ora bakal ngasilake suku lambda murni. Mula cara iki
ora bisa digunakake kanggo netepake faktorial, utawa
fungsi rekursif umume.

Nanging definisi rekursif mau isih menehi sawijining pituduh:
Yen $f$ suku kang makili fungsi faktorial, suku
\[
\fn{Fac}' \ident \lambd[g][\lambd[n][\fn{IsZero} \, n \, \num 1\, (\fn{Mult} \, n\, (g (\fn{Pred} n)))]]
\]
kang ditrapake marang suku $f$, yaiku $\fn{Fac}'\,f$, uga makili
fungsi faktorial. Tegese, yen $\fn{Fac}'$ dianggep minangka fungsi
kang nampa fungsi lan ngasilake fungsi, nilaine
$\fn{Fac'}\, f$ padha karo~$f$, yen $f$ iku fungsi faktorial.
Fungsi~$f$ kang nduweni sipat $\fn{Fac}' \, f \equal[\beta] f$ diarani
\emph{titik tetep} saka $\fn{Fac}'$. Dadi, fungsi faktorial iku
titik tetep saka $\fn{Fac}'$.

Ing kalkulus lambda ana suku-suku kang ngitung titik tetep
saka suku kang diwenehake. Suku-suku iki bisa digunakake
supaya suku kaya $\fn{Fac}'$ dadi definisi faktorial.

\begin{defn}
  \ollabel{defn:Turing-Y}
  \emph{Kombinator-Y} yaiku suku iki:
  \[
  Y \ident (\lambd[ux][x(uux)])(\lambd[ux][x(uux)]).
  \]
\end{defn}

\begin{thm}
  $Y$ nduweni sipat $Yg \red g(Yg)$ kanggo saben suku $g$.
  Mula $Yg$ tansah dadi titik tetep saka~$g$.
\end{thm}

\begin{proof}
  Cekakna $(\lambd[ux][x(uux)])$ dadi~$U$, saéngga $Y \ident UU$. Banjur
  \begin{align*}
    Y g &\ident (\lambd[ux][x(uux)])U\,g \\
    &\red (\lambd[x][x(UUx)])g\\
    & \red g(UUg) \ident g(Yg).
  \end{align*}
  Amarga $g(Yg)$ lan $Yg$ loro-lorone bisa direduksi dadi $g(Yg)$,
  mula $g(Yg) \equal[\beta] Yg$; dadi $Yg$ titik tetep saka~$g$.
\end{proof}

Mesthi wae, amarga $Yg$ ngemot redex, reduksi bisa terus tanpa wates:
\begin{align*}
  Y g &\red g (Y g) \\
    &\red g (g (Y g)) \\
    &\red g(g (g (Y g)))\\
    &\ldots
\end{align*}
Mula $Yg$ bisa dibayangake minangka $g$ kang ditrapake marang
awake dhewe tanpa wates kaping. Yen $g$ ditrapake sepisan maneh,
kaya-kaya ora ana tambahan apa-apa: $g$ kang ditrapake marang~$g$
kang wis ditrapake tanpa wates kaping marang~$Yg$ isih padha karo
$g$ kang ditrapake tanpa wates kaping marang~$Yg$.

Gatekna yen rerangkening langkah reduksi-$\beta$ ing ndhuwur kang
diwiwiti saka~$Yg$ iku tanpa wates. Dadi yen $Yg$ ditrapake marang
sawijining suku, yaiku $(Yg)N$, suku iku uga bisa direduksi dadi
tanpa wates suku kang beda-beda, yaiku $(g(Yg))N$,
$(g(g(Yg)))N$, \dots. Nanging bisa wae ana rerangkening reduksi
\emph{liyane} kang mandheg ing wujud normal.

Jupuken faktorial minangka conto. Tetepna $\fn{Fac}$ minangka
$Y\, \fn{Fac}'$ (yaiku titik tetep saka $\fn{Fac'}$). Banjur:
\begin{align*}
  \fn{Fac}\,\num 3
  &\red Y\, \fn{Fac}' \, \num 3 \\
  &\red \fn{Fac}' (Y \,\fn{Fac}') \, \num 3 \\
  &\ident (\lambd[x][\lambd[n][\fn{IsZero} \, n\, \num 1\,
      (\fn{Mult}\, n\, (x (\fn{Pred}\, n)))]]) \, \fn{Fac} \, \num 3\\
  & \red \fn{IsZero} \, \num 3\, \num 1\,
      (\fn{Mult}\, \num 3\, (\fn{Fac} (\fn{Pred}\, \num 3))) \\
  &\red  \fn{Mult}\, \num 3 \, (\fn{Fac} \, \num 2).
  \intertext{Kanthi cara kang padha,}
  \fn{Fac}\,\num 2
  &\red  \fn{Mult}\, \num 2 \, (\fn{Fac} \, \num 1) \\
  \fn{Fac}\,\num 1
  &\red  \fn{Mult}\, \num 1 \, (\fn{Fac} \, \num 0)
  \intertext{nanging}
  \fn{Fac}\,\num 0
  &\red \fn{Fac}' (Y \,\fn{Fac}') \, \num 0 \\
  &\ident (\lambd[x][\lambd[n][\fn{IsZero} \, n\, \num 1\,
      (\fn{Mult}\, n\, (x (\fn{Pred}\, n)))]]) \, \fn{Fac} \, \num 0\\
  & \red \fn{IsZero} \, \num 0\, \num 1\,
      (\fn{Mult}\, \num 0\, (\fn{Fac} (\fn{Pred}\, \num 0))).\\
  &\red \num 1.
  \intertext{Mula yen digabungake,}
  \fn{Fac}\,\num 3
  &\red  \fn{Mult}\, \num 3 \,
  (\fn{Mult} \, \num 2\, (\fn{Mult} \, \num 1\, \num 1)).
\end{align*}

Bab kang lumaku kanggo $\fn{Fac'}$ uga lumaku kanggo saben
definisi rekursif. Upamakna kita duwe persamaan rekursif
\begin{align*}
g\,x_1\dots x_n & \equal[\beta] N
\intertext{kanthi $N$ bisa ngemot~$g$ lan $x_1$, \dots,~$x_n$. Mula mesthi
ana suku~$G \ident (Y \lambd[g][\lambd[x_1\dots x_n][N]])$ kang
nyukupi}
G\,x_1 \dots x_n & \equal[\beta] \Subst{N}{G}{g}.
\intertext{Amarga miturut teorema titik tetep,}
G \ident (Y \lambd[g][\lambd[x_1\dots x_n][N]]) & \red \lambd[g][\lambd[x_1\dots x_n][N]](Y \lambd[g][\lambd[x_1\dots x_n][N]])\\
& \ident (\lambd[g][\lambd[x_1\dots x_n][N]])\,G
\intertext{lan mulane}
G\,x_1 \dots x_n 
& \red (\lambd[g][\lambd[x_1\dots x_n][N]])\,G\,x_1\dots x_n\\
& \red (\lambd[x_1\dots x_n][\Subst{N}{G}{g}])\,x_1\dots x_n\\
  & \red \Subst{N}{G}{g}.
\end{align*}

Kombinator $Y$ ing \olref{defn:Turing-Y} asalé saka Alan
Turing. Alonzo Church ngajokake versi liya kang kita
arani~$Y_C$:
\[
Y_C \ident \lambd[g][(\lambd[x][g(xx)])(\lambd[x][g(xx)])].
\]
Kombinator Church rada luwih ringkih tinimbang kombinator Turing:
$Y_Cg \equal[\beta] g(Y_Cg)$, nanging ora $Y_Cg \bred g(Y_Cg)$.
Upamakna $V$ iku suku $\lambd[x][g(xx)]$, saéngga
$Y_C \ident \lambd[g][VV]$, kanthi cekakan iku dimekarake
ing sangisore pengiket paramètere. Banjur
\begin{align*}
  VV & \ident (\lambd[x][g(xx)])V \red g(VV)
  \text{ lan mulane}\\
  Y_C g & \ident (\lambd[g][VV])g \red VV \red g(VV), \text{ nanging uga}\\
  g(Y_C g) & \ident g((\lambd[g][VV])g) \red g(VV).
\end{align*}
Tegese, $Y_Cg$ lan $g(Y_Cg)$ loro-lorone direduksi dadi
suku kang padha, yaiku $g(VV)$; mula $Y_Cg \equal[\beta] g(Y_Cg)$.
Sipat iki asring cukup kanggo panganggone.

\end{document}
